\documentclass[10pt]{article}
\usepackage[margin=1in]{geometry}
\usepackage[T1]{fontenc}
\usepackage{lmodern,microtype,amsmath,amssymb,amsthm}
\theoremstyle{definition}
\newtheorem{assumption}{Assumption}
\begin{document}
\section*{Stochastic Constrained Quasar Minimization}
We consider minimizing $f:E\to\mathbb{R}$ over a compact convex set $X$ in a finite-dimensional real Hilbert space $E$. On $X$, the function has an $L$-Lipschitz gradient, where $L>0$, and is $\gamma$-quasar-convex relative to a minimizer $x^\star\in X$, where $0<\gamma\leq1$. The oracles return unbiased stochastic gradients with bounded variance using independent, identically distributed samples, and exact Euclidean projections onto $X$. Given $\varepsilon>0$ and $\delta\in(0,1)$, the goal is a measurable output $\bar{x}$ satisfying both $\bar{x}\in X$ and $f(\bar{x})-f(x^\star)\leq\varepsilon$ with probability at least $1-\delta$ over sampling and algorithm randomness.

The optimization problem is
\[
\min_{x\in X} f(x).
\]
Write $\langle\cdot,\cdot\rangle$ for the inner product on $E$ and $\|\cdot\|$ for its induced norm. The algorithm receives $E$, the measurable sample space $\Omega$, a starting point $x^0$, and the parameters $L,D>0$, $\gamma\in(0,1]$, $\sigma\geq0$, $\varepsilon>0$, and $\delta\in(0,1)$; the objective, feasible set, minimizer, and sampling distribution are accessed through the oracles.

Our assumptions:
\begin{assumption}[Feasible geometry]
The set $X\subseteq E$ is compact and convex, contains $x^0$ and $x^\star$, and satisfies
\[
\|x-y\|\leq D,\qquad x,y\in X.
\]
\end{assumption}

\begin{assumption}[Smoothness and quasar convexity]
The function $f$ is Fr\'echet differentiable at every point of $X$, and, for every $x,y\in X$,
\[
\begin{aligned}
\|\nabla f(x)-\nabla f(y)\|&\leq L\|x-y\|,\\
f(x^\star)&\leq f(x),\\
\gamma\bigl(f(x)-f(x^\star)\bigr)
&\leq\langle\nabla f(x),x-x^\star\rangle.
\end{aligned}
\]
\end{assumption}

\begin{assumption}[Stochastic gradient oracle]
There are a probability measure $\nu$ on $\Omega$ and a jointly measurable map $g:E\times\Omega\to E$. For every $x\in X$, the random vector $g(x,\xi)$ is Bochner integrable and $\|g(x,\xi)-\nabla f(x)\|^2$ is integrable under $\xi\sim\nu$, with
\[
\begin{aligned}
\mathbb{E}_{\xi\sim\nu}[g(x,\xi)]&=\nabla f(x),\\
\mathbb{E}_{\xi\sim\nu}
 [\|g(x,\xi)-\nabla f(x)\|^2]&\leq\sigma^2.
\end{aligned}
\]
The oracle uses a sequence $(\xi_k)_{k\geq0}$ of measurable, mutually independent sample variables, each with law $\nu$, independent of the algorithm's randomness. At a stochastic gradient call with query $x$, it returns $g(x,\xi_k)$ using the sample at that call's index.
\end{assumption}

\begin{assumption}[Exact projection oracle]
The projection oracle returns $\Pi_X(z)$ at $z\in E$, where
\[
\begin{aligned}
\Pi_X(z)&\in X,\\
\langle z-\Pi_X(z),w-\Pi_X(z)\rangle&\leq0,
\qquad w\in X.
\end{aligned}
\]
\end{assumption}

Both oracles accept queries throughout $E$; the gradient and moment conditions above apply on $X$. Stochastic gradient calls and projection calls have separate budgets. After these calls, the output $\bar{x}\in E$ must be measurable and satisfy
\[
P\bigl\{\bar{x}\notin X\ \text{or}\
 f(\bar{x})-f(x^\star)>\varepsilon\bigr\}\leq\delta,
\]
where $P$ is the joint law of the samples and the independent algorithm randomness.
\end{document}
